\paragraph{Capacity information versus cohort selection.}
We recomputed minimum fixed budgets on the same compatible records, retaining
the uncapped width target. Table~\ref{tab:paired-capacity} separates excluded
members of the original 225 reading-saving blocks from those for which the
cap makes a fixed batch sufficient. At $E/2$, only 7 original and 2 external
members still require feedback. The full capped counts of 56 and 4 also
include 49 and 2 previously unguaranteed blocks that become adaptive-certifiable.
At $E/3$, every compatible member of the original 225 becomes fixed-certifiable.
The minimum fixed totals on identical compatible subsets fall from 2,509 to
1,840 and from 543 to 424 at $E/2$, and from 2,003 to 768 and from 538 to 220
at $E/3$, for the original and external panels respectively. This attributes
the change to both exclusions and genuinely stronger information.

\begin{table}[htbp]\centering\small
\caption{Capacity transitions within the original 225 reading-saving blocks.
The last two columns partition compatible members. The cap and precision
target are the same for both acquisition policies.}
\label{tab:paired-capacity}
\begin{tabular}{llrrrr}\toprule
Panel & Cap & Reference & Excluded & Fixed $\leq2$ & Still needs feedback\\\midrule
Original & $E/2$ & 168 & 8 & 153 & 7\\
External & $E/2$ & 57 & 17 & 38 & 2\\
Original & $E/3$ & 168 & 58 & 110 & 0\\
External & $E/3$ & 57 & 17 & 40 & 0\\\bottomrule
\end{tabular}
\end{table}
